\section{Discussion and Conclusion}
\label{sec:discussion}

% Paragraph 1: Summary of the Core Finding & Primary Method
% - Restate the research question (e.g., "Does context matter for optimization?").
% - Summarize the key insight (e.g., "We found that signal patterns vary systematically by task, explaining why static baselines fail.").
% - State that Method A successfully exploits this insight to outperform baselines and avoid specific failure modes (cite specific numbers/deltas here).
We explored \textbf{[Core Question]}, finding that \textbf{[Phenomenon]} varies systematically across domains. \textbf{[Method Name]} addresses this by leveraging \textbf{[Key Signal]}, consistent preserving performance within \textbf{[Delta]} of the original model while outperforming \textbf{[Baselines]} and avoiding \textbf{[Specific Failure Mode]}.

% Paragraph 2: Theoretical Validation (The Oracle/Variant)
% - Discuss the secondary variant/oracle (Method B).
% - Explain that this variant serves as an upper bound or proof-of-concept.
% - Mention if it reveals something deeper about the problem (e.g., "It confirms that sensitivity is the decisive factor"), proving the theoretical premise even if the variant isn't for deployment.
Additionally, \textbf{[Oracle Variant]} establishes an upper bound, confirming that \textbf{[Core Metric]} is decisive. While primarily for validation, this result proves that \textbf{[Optimization Process]} can serve as a lens to understand how \textbf{[Task Information]} is encoded, validating the premise for lightweight proxies.

% Paragraph 3: Broader Implications, Limitations, and Future Work
% - Highlight a side benefit (e.g., Interpretability: "The method also acts as a probe").
% - Acknowledge 1-2 limitations (e.g., "Dependence on calibration data").
% - Suggest 1-2 future directions (e.g., "Dynamic allocation").
% - Final concluding sentence.
Beyond efficiency, this approach offers \textbf{[Secondary Benefit/Interpretability]}. Limitations include \textbf{[Limitation A]} and \textbf{[Limitation B]}. Promising future directions include \textbf{[Future Work A]}. Overall, our work demonstrates that \textbf{[Context-Awareness]} is essential for preserving capabilities where it matters most.